\section{One field prescription for every truncation degree}
\label{sec:ope-prescription}
An operator-product expansion must use the same coefficients when its
truncation degree changes. Choosing a new basis independently at each
finite degree would not ensure this. We construct a single compatible
system of fields and finite correlation tests, one degree at a time.
The tests are fixed before choosing the cluster geometry or its spectators.

Let $E_n$ be the finite-dimensional complex space of physical fields through
degree $n$. Denote its injective inclusion maps by
$\iota_{mn}:E_m\to E_n$, $m\leq n$, with
$\iota_{nn}=I$ and $\iota_{nk}\iota_{mn}=\iota_{mk}$.
These are the physical filtration maps, so inserting a field and its image
in a higher degree gives the same separated correlations.

An exterior test $\lambda$ consists of finitely many fixed spectator labels
and a separated configuration including a distinguished first position.
Inserting $v\in E_n$ at that first position and evaluating the mixed
Euclidean correlation defines a linear functional $\lambda_n(v)$.
The physical field quotient has been defined by vanishing of all such
correlations. Thus these functionals separate $E_n$:
\begin{equation}\label{eq:ope-detector-separation}
 [\lambda_n(v)=0\text{ for every exterior test }\lambda]\ \Longrightarrow\ v=0.
\end{equation}
Their compatibility with filtration is
$\lambda_n\iota_{mn}=\lambda_m$. For a finite linear combination
$b=\sum_{\lambda\in F}b_\lambda\lambda$ use the notation
$b_n=\sum_{\lambda\in F}b_\lambda\lambda_n$. The same finite combination
then defines a test on every level.

\begin{lemma}\label{lem:ope-detector-span}
Every linear functional on $E_n$ is $b_n$ for some finite combination $b$
of exterior tests.
\end{lemma}
\begin{proof}
Let $V$ be the span of the exterior functionals inside $E_n^*$.
If $V$ were proper, a basis of $V$ would consist of fewer than $\dim E_n$
functionals. The linear map from $E_n$ to their coordinate values would
then have a nonzero kernel, by rank--nullity. A nonzero vector in that
kernel would be annihilated by all of $V$, hence by every exterior test,
contrary to \eqref{eq:ope-detector-separation}. Therefore $V=E_n^*$.
Membership in a span means a finite linear combination, as required.
\end{proof}

Put $E_{-1}=0$ and
$d_n=\dim(E_n/\iota_{n-1,n}E_{n-1})$, where the denominator is zero
when $n=0$. A channel is a pair $\alpha=(n,j)$ with $1\leq j\leq d_n$;
its degree is $|\alpha|=n$. We seek fields $u_\alpha\in E_{|\alpha|}$
and finite exterior tests $b_\alpha$ such that, for every $n$,
\begin{equation}\label{eq:ope-biorthogonal}
 b_{\alpha,n}(\iota_{|\beta|,n}u_\beta)=\delta_{\alpha\beta}
       \quad(|\alpha|,|\beta|\leq n),
\end{equation}
and
\begin{equation}\label{eq:ope-rebuild}
 v=\sum_{|\alpha|\leq n}b_{\alpha,n}(v)
                       \iota_{|\alpha|,n}u_\alpha\qquad(v\in E_n).
\end{equation}
The sums contain finitely many channels. No assertion of a convergent
infinite expansion is involved.

\begin{proposition}\label{prop:ope-prescription}
There is one choice of all the fields $u_\alpha$ and tests $b_\alpha$
satisfying \eqref{eq:ope-biorthogonal} and \eqref{eq:ope-rebuild}
simultaneously for every $n$. Moreover
$b_{\alpha,m}=0$ whenever $m<|\alpha|$.
\end{proposition}
\begin{proof}
At degree zero choose a basis of $E_0$ and represent each coordinate
functional by a finite exterior test, using
Lemma~\ref{lem:ope-detector-span}. The two identities are then the usual
basis identities.

Suppose the fields and tests through degree $n$ have been chosen.
Write $R=\iota_{n,n+1}$, let $q:E_{n+1}\to E_{n+1}/R(E_n)$ be the
quotient map, and set
\[
 Bv=\sum_{|\alpha|\leq n}b_{\alpha,n+1}(v)
                           \iota_{|\alpha|,n}u_\alpha.
\]
Test compatibility and \eqref{eq:ope-rebuild} give $BR=I_{E_n}$.
Choose a linear right inverse $S$ of $q$; it exists by lifting a finite
basis of the quotient. Correct this arbitrary lift by putting
\[
 F=S-RBS.
\]
Then $qF=I$ and $BF=0$. The old tests vanish on the new image $F$.
Indeed \eqref{eq:ope-biorthogonal} gives
$b_{\alpha,n}B=b_{\alpha,n+1}$ for every old channel $\alpha$, so
$b_{\alpha,n+1}F=b_{\alpha,n}BF=0$.

Every $v\in E_{n+1}$ has the exact decomposition
\begin{equation}\label{eq:ope-quotient-split}
 v=Fqv+RBv.
\end{equation}
To prove it, $q(v-Fqv)=0$ implies $v-Fqv=Ra$ for some $a\in E_n$.
Applying $B$ and using $BF=0$, $BR=I$ gives $a=Bv$.
Now choose a basis $(e_j)$ of the quotient. The new fields are $Fe_j$.
For each new channel choose a finite exterior test representing
$e_j^*q$ on $E_{n+1}$. Leave every old field and every old test unchanged.
The new tests annihilate $R(E_n)$, are dual to the new fields because
$qF=I$, and the old tests annihilate the new fields as just proved.
Old fields and tests retain their mutual duality by compatibility.
This proves all four old/new cases of \eqref{eq:ope-biorthogonal}.
Expanding the two terms in \eqref{eq:ope-quotient-split} proves
\eqref{eq:ope-rebuild} at degree $n+1$.

Iterate. A channel and its finite test, once chosen, never change.
Their value in the resulting infinite list is the one fixed at their birth
degree. Any identity at degree $n$ involves only the finitely many channels
already fixed by stage $n$, so it remains valid in the final list.
Finally take $m<|\alpha|$ and expand $v\in E_m$ using
\eqref{eq:ope-rebuild}. Lift that identity to degree $|\alpha|$ and apply
$b_\alpha$. Biorthogonality makes every term zero, because all its
channels have degree at most $m$. Compatibility then gives
$b_{\alpha,m}(v)=0$.
\end{proof}
The subtraction $RBS$ is the reason this construction works at all degrees.
It places every new quotient lift in the exact kernels of all old tests,
while retaining the old tests as actual finite correlation evaluations.

\subsection{Testing a product at the cluster center}
For each exterior test $\lambda$, translate its distinguished position to
zero. Write its other positions as $z_{\lambda,j}$ and its spectator
labels as $w_{\lambda,j}$. These positions are pairwise distinct and none
is zero. For a cluster of fields $v_i$ at $x_i$ and center $o$, define
\[
 \mathcal R_\lambda(\boldsymbol x,\boldsymbol v;o)
 =S_{\mathrm{mix}}\bigl((v_i,x_i-o)_i,
                         (w_{\lambda,j},z_{\lambda,j})_j\bigr)
\]
when the combined configuration is separated, and assign value zero
otherwise. This convention gives a function everywhere; every asymptotic
use below is in its separated region. Define
\begin{equation}\label{eq:ope-raw-coefficient}
 c_\alpha^{\mathrm{raw}}(\boldsymbol x,\boldsymbol v;o)
       =\sum_{\lambda\in\operatorname{supp}b_\alpha}
                  b_{\alpha,\lambda}
                  \mathcal R_\lambda(\boldsymbol x,\boldsymbol v;o),
\end{equation}
and the finite field-valued expression
\begin{equation}\label{eq:ope-raw-shadow}
 H_n(\boldsymbol x,\boldsymbol v;o)
   =\sum_{|\alpha|\leq n}
       c_\alpha^{\mathrm{raw}}(\boldsymbol x,\boldsymbol v;o)
                                      \iota_{|\alpha|,n}u_\alpha.
\end{equation}
Each coefficient is fixed independently of the truncation degree, the
phase, or any later choice of exterior spectators. It is multilinear in
the cluster labels on its separated region. The separation condition
depends only on positions, so the zero extension is multilinear as well.
Permutation symmetry of the correlations gives simultaneous permutation
symmetry in cluster positions and labels. Simultaneous translation of the
cluster and center leaves their differences unchanged.

For a single field $v\in E_m$ placed at $o$, translation invariance and
the definition of the exterior test give
\[
 c_\alpha^{\mathrm{raw}}((o),(v);o)=b_{\alpha,m}(v).
\]
All the detector configurations in this evaluation are separated. Hence,
for $m\leq n$, \eqref{eq:ope-rebuild} implies
\begin{equation}\label{eq:ope-raw-single}
 H_n((o),(v);o)=\iota_{mn}v.
\end{equation}
The same testing operation is compatible with replacing any input field
by its filtration image, since separated correlations have that
compatibility.

For every fixed finite list of coefficients there is a common positive
detector clearance. To see this, take the minimum of the finitely many
nonzero lengths $|z_{\lambda,j}|$ occurring in their finite test supports;
if there are no spectator positions, there is no such restriction.
If $|x_i-o|$ is less than half that minimum, no cluster point can meet a
detector position. Distinctness among cluster points and among detector
points supplies the remaining separation conditions. The same argument
works after any orthogonal rotation of the relative cluster, because
rotations preserve $|x_i-o|$. Smoothness of the mixed correlations on
separated configurations makes each such finite coefficient collection
smooth on this common region. This is a finite-degree statement; it gives
no uniform clearance for infinitely many degrees.

The raw expressions above do not yet have the required rotational and
reality properties, because the finite tests were chosen by linear
algebra. Averaging the field-valued expression over $O(4)$ and symmetrizing
conjugation supplies those properties. The subsequent homogeneous
projection must be made using these averaged tests themselves, so that
its lower-degree components agree exactly. These operations and their
quantitative error estimates are the next steps of the OPE construction.
